% GENERATED FILE. Edit canonical lessons, then run npm run generate:books.
% canonical-source-hash: fnv1a64:c51fe60ae8424290
% canonical-lessons: MR-W64-ddha, MR-W64-pha, MR-W64-danda, MR-R64-rows-closed

\chapter{Rows Closed}
\label{ch:mr-oli-purna}
\hlchaptermodality{64}

\begin{chapteropening}
\textbf{By the end of this chapter:} \emph{I can write the two aspirated letters that were missing from the curled and lip rows, and recognise the danda where an older text ends a sentence.}

From cold, with nothing open: write the aspirated retroflex d, the aspirated lip p, and the mark an older text ends a sentence with.
\end{chapteropening}

\section[ddha]{\mr{ढ} — the fifth of the curled row}

\label{lesson:MR-W64-ddha}



\begin{quote}
Chapter seven gave you \textbf{\mr{ट}}, \textbf{\mr{ठ}}, \textbf{\mr{ड}} and \textbf{\mr{ण}}. Four of a row of five. Say the two of those four that belong to the \emph{d} half of the row.
\end{quote}

\subsection*{Script}
\begin{quote}
\mr{ढ}
\end{quote}

\textbf{\mr{ढ}} is the aspirated partner of \textbf{\mr{ड}} — the same step you have taken three times already, at \textbf{\mr{क}}/\textbf{\mr{ख}}, at \textbf{\mr{ट}}/\textbf{\mr{ठ}} and at \textbf{\mr{त}}/\textbf{\mr{थ}}. Curl the tongue to the roof of the mouth for a \textbf{\mr{ड}}, then let the breath out after it.

Set the two side by side and the shape follows the sound: \textbf{\mr{ड} \mr{ढ}}. The second carries an extra curl at its foot, so the letter that has more breath in it also has more ink.

\textbf{\mr{ढग}} — \emph{a cloud}. Two letters, and you can read both.

\begin{grammarlens}[title={Grammar Lens: the romanization lies about this one}]
\noindent
\begin{tabularx}{\linewidth}{@{}>{\raggedright\arraybackslash}X>{\raggedright\arraybackslash}X@{}}
\toprule
\textbf{letter} & \textbf{where the tongue is} \\
\midrule
\mr{ध} & against the teeth \\
\mr{ढ} & curled to the roof \\
\bottomrule
\end{tabularx}

Both of these are written \textbf{dh} in most romanizations, and that is the whole of the problem. \textbf{They are different letters in different rows}, and the spelling that is supposed to help you tell them apart makes them identical.

\textbf{On the page they are nothing alike}, which is the good news: a reader who learns the shapes is never in doubt, and a reader who leans on the romanization is in doubt every time. This is the argument for the script the track has been making since chapter one, in one pair of letters.
\end{grammarlens}

\subsection*{Writing: guided copy}
Copy \textbf{\mr{ढ}} twice with the model visible. Then write \textbf{\mr{ड}} beside it and check that only one of the two carries the extra curl at the foot.

\begin{tcolorbox}[breakable,colback=teal!4,colframe=teal!35!black,title={Before you move on}]
Which of the two has more ink at the foot? (\textbf{\mr{ढ}}.) Which of \textbf{\mr{ध}} and \textbf{\mr{ढ}} is made with the tongue against the teeth? (\textbf{\mr{ध}}.) What do most romanizations write for both of them? (\emph{dh} — which is why the shapes matter.)

Stroke order for \textbf{\mr{ढ}} is illustrated in the Devanagari letter-order animations held by Wikimedia Commons.
\end{tcolorbox}
\section[pha]{\mr{फ} — the gap in the lip row}

\label{lesson:MR-W64-pha}



\begin{quote}
The lip row is \textbf{\mr{प}}, \textbf{\mr{फ}}, \textbf{\mr{ब}}, \textbf{\mr{भ}}, \textbf{\mr{म}}. Name the three of those five you have been writing since the early chapters. Then say what the lesson before this one did to \textbf{\mr{ड}}.
\end{quote}

\subsection*{Script}
\begin{quote}
\mr{फ}
\end{quote}

\textbf{\mr{फ}} is \textbf{\mr{प}} with the breath let out. That is the fifth time this same step has appeared, and with it every row of stops in the script is closed:

\noindent
\begin{tabularx}{\linewidth}{@{}>{\raggedright\arraybackslash}X>{\raggedright\arraybackslash}X@{}}
\toprule
\textbf{plain} & \textbf{aspirated} \\
\midrule
\mr{क} \mr{ट} \mr{त} & \mr{ख} \mr{ठ} \mr{थ} \\
\mr{ड} \mr{प} & \mr{ढ} \mr{फ} \\
\bottomrule
\end{tabularx}

Five pairs, one rule, and it has not changed once: \textbf{the aspirated letter is the plain one with a puff of air after it.} A row is not a list of five unrelated shapes — it is two sounds crossed with two settings, plus the nasal at the end.

\begin{grammarlens}[title={Grammar Lens: the letter whose sound has drifted}]
\textbf{\mr{फोन}} — \emph{a telephone}. Read the first letter aloud and listen to what you did with your lips.

In words Marathi has borrowed, many speakers give \textbf{\mr{फ}} the sound of an English \emph{f} rather than the puffed \emph{p} the row predicts. Both are heard, and the drift runs one way: it is the borrowed words that pull toward \emph{f}, while the inherited ones keep the puff.

\textbf{The letter did not change; the words around it did.} A reader meets one shape and has to know that it covers two sounds, which is the opposite of the tidy row above and worth expecting rather than being surprised by.
\end{grammarlens}

\subsection*{Writing: guided copy}
Copy \textbf{\mr{फ}} twice with the model visible. Then write \textbf{\mr{प}} beside it and read the two aloud, plain then aspirated.

\begin{tcolorbox}[breakable,colback=teal!4,colframe=teal!35!black,title={Before you move on}]
Which letter is \textbf{\mr{प}} with a puff after it? (\textbf{\mr{फ}}.) How many plain and aspirated pairs does the script have now? (Five.) What sound does \textbf{\mr{फ}} take in a borrowed word like \textbf{\mr{फोन}}? (Often an \emph{f}.)

Stroke order for \textbf{\mr{फ}} is illustrated in the Devanagari letter-order animations held by Wikimedia Commons.
\end{tcolorbox}
\section[danda]{\mr{।} — the stroke that ends an older sentence}

\label{lesson:MR-W64-danda}



\begin{quote}
Name the two marks you have been writing that Marathi borrowed whole from the Latin alphabet — the one that asks and the one that lists.
\end{quote}

\subsection*{Script}
\begin{quote}
\mr{।}
\end{quote}

\textbf{\mr{।}} is the \textbf{danda} — a single upright stroke, drawn from the headline down, and it ends a sentence. It is the full stop Devanagari had before there was a full stop.

Everything printed in this track ends its sentences with a \textbf{.}, because that is what Marathi does now. \textbf{The danda did not go away; it moved.} It stays in poetry, in older printing, in devotional and formal text, and on plenty of signs and notices that were set once and never reset.

\begin{grammarlens}[title={Grammar Lens: a mark to recognise rather than to reach for}]
\noindent
\begin{tabularx}{\linewidth}{@{}>{\raggedright\arraybackslash}X>{\raggedright\arraybackslash}X>{\raggedright\arraybackslash}X@{}}
\toprule
\textbf{mark} & \textbf{what it does} & \textbf{do you write it} \\
\midrule
. & ends a sentence & yes, every time \\
\mr{।} & ends a sentence & you will read it \\
\bottomrule
\end{tabularx}

\textbf{This is the first mark in the track you are being taught to read and not to produce}, and the distinction is worth naming because it will come back. A reader who has only ever met the full stop hits a danda on the first authentic notice and has no way of knowing whether the sentence ended, whether a word is missing, or whether something has gone wrong with the printing.

It is not a vertical bar between columns and it is not the tail of a letter. It sits alone after the last word, with a space before whatever comes next, and it means what a full stop means.
\end{grammarlens}

\subsection*{Writing: guided copy}
Copy \textbf{\mr{।}} twice with the model visible. Draw it from the headline straight down, the full height of a letter, with no hook at either end. Then write a comma beside it and check that one sits on the baseline and the other runs the whole way down.

\begin{tcolorbox}[breakable,colback=teal!4,colframe=teal!35!black,title={Before you move on}]
What does a danda mean? (The sentence ended.) Which of the two sentence-enders will you be writing? (The full stop.) Where are you most likely to meet the other one? (Poetry, older printing, and signs nobody has reset.)
\end{tcolorbox}
\section[olī pūrṇa]{rows closed — two letters and a stroke, from cold}

\label{lesson:MR-R64-rows-closed}



\begin{quote}
Close the three lessons before this one. Write the plain retroflex \textbf{d} and the plain lip \textbf{p} from memory, with nothing in front of you.
\end{quote}

\subsection*{Script: the five pairs, written out once}
\noindent
\begin{tabularx}{\linewidth}{@{}>{\raggedright\arraybackslash}X>{\raggedright\arraybackslash}X@{}}
\toprule
\textbf{plain} & \textbf{aspirated} \\
\midrule
\mr{क} \mr{ट} \mr{त} & \mr{ख} \mr{ठ} \mr{थ} \\
\mr{ड} \mr{प} & \mr{ढ} \mr{फ} \\
\bottomrule
\end{tabularx}

Ten letters, one rule, and the last two of them arrived in this chapter. \textbf{Every stop in the script now has its pair}, which means a plain letter you meet from here on has a partner you can predict rather than a shape you have to learn.

Write the bottom row from memory, plain beside aspirated, and check that the aspirated one in each pair is the one carrying the extra ink.

\begin{grammarlens}[title={Grammar Lens: the pair the romanization hides}]
\textbf{\mr{ध}} and \textbf{\mr{ढ}} are both \emph{dh} in a romanization and are not the same letter. One is made against the teeth and one with the tongue curled to the roof, and the page tells them apart where the spelling will not.

Say \textbf{\mr{ढग}} and then say a word with \textbf{\mr{ध}} in it, and listen for where your tongue went.
\end{grammarlens}

\subsection*{Your turn}
Say these aloud:

\begin{itemize}
\raggedright
  \item which of \textbf{\mr{ध}} and \textbf{\mr{ढ}} is the curled one
  \item what sound \textbf{\mr{फ}} often takes in a borrowed word
  \item what a \textbf{\mr{।}} at the end of a line means
  \item which of the two sentence-enders you write yourself
\end{itemize}

\begin{tcolorbox}[breakable,colback=teal!4,colframe=teal!35!black,title={Before you move on}]
How many plain-and-aspirated pairs does the script have? (Five.) Which mark in this chapter are you learning to read rather than to write? (The \textbf{\mr{।}}.)
\end{tcolorbox}
